%%%PAGE 240 rot=0 legibility=good summary=共点力平衡专题：活结与死结、平衡条件与方法（剪断绳子、三力/多力、相似三角形、摩擦角4→3）；圆锥摆与双绳圆周运动的ω范围
%%%BLOCK id=240-01 ch=02 sec=共点力平衡·活结与死结 kind=knowledge alt=none cont=none y=0.01-0.15
% 左上角淡铅笔小字标题
\textbf{活结}

%FIG p240_a 0.06 0.03 0.17 0.108
\notefig[0.2]{figures/p240_a.png}

活结两侧拉力相等

%FIG p240_b 0.222 0.046 0.318 0.102
\notefig[0.2]{figures/p240_b.png}

$T\cos\theta=\frac12 mg$　　$T=\dfrac{mg}{2\cos\theta}$　　$2T\cos\theta=mg$　　$d\uparrow\ \cos\alpha\downarrow\ T\uparrow$ % CHECK: 末式中的角手写像 α（前式为 θ），照录

$\sin\theta=\dfrac{d}{L}$

%FIG p240_c 0.555 0.108 0.645 0.15
\notefig[0.2]{figures/p240_c.png}

引体向上　\unclear{竖}手力小
%%%END
%%%BLOCK id=240-02 ch=02 sec=共点力平衡·活结与死结 kind=knowledge alt=none cont=none y=0.16-0.25
%FIG p240_d 0.075 0.158 0.205 0.25
\notefig[0.25]{figures/p240_d.png}

$d$先$\uparrow$后$\downarrow$

$T$先$\uparrow$后$\downarrow$
%%%END
%%%BLOCK id=240-03 ch=02 sec=共点力平衡·活结与死结 kind=knowledge alt=none cont=none y=0.25-0.37
死结　结不沿绳移动，为两根绳

活杆　沿杆

死杆　任意

%FIG p240_e 0.52 0.255 0.68 0.368
\notefig[0.2]{figures/p240_e.png}
%%%END
%%%BLOCK id=240-04 ch=02 sec=共点力平衡·条件与方法 kind=knowledge alt=none cont=none y=0.38-0.46
共点力（延长线交于1点）平衡 $\Leftrightarrow$ 缓慢　匀速、静止、$F_{\text{合}}=0$ % 原稿"1点"写在括号内第二行

方法：二力平衡条件　等效
%%%END
%%%BLOCK id=240-05 ch=02 sec=共点力平衡·条件与方法 kind=method alt=03 cont=none y=0.47-0.57
%FIG p240_f 0.10 0.488 0.245 0.568
\notefig[0.2]{figures/p240_f.png}

平衡后剪断绳子，$F_B=mg$

%FIG p240_g 0.575 0.462 0.835 0.55
\notefig[0.3]{figures/p240_g.png}
% 斜面简图旁手写：α t_1, t_2 α（斜面上方），t_1（块旁）；弯箭头自"验证"处指向左侧 F_B=mg
$\alpha\ \checkmark$　验证
%%%END
%%%BLOCK id=240-06 ch=02 sec=共点力平衡·条件与方法 kind=method alt=none cont=none y=0.57-0.66
三力平衡　力三角形

多力平衡　正交分解

相似，力辅助圆，正弦定理、余弦定理　勾股定理　　$\sin(\pi-\alpha)=\sin\alpha$
%%%END
%%%BLOCK id=240-07 ch=02 sec=共点力平衡·摩擦角 kind=derivation alt=none cont=none y=0.67-0.78
$4\to3$

%FIG p240_h 0.14 0.683 0.44 0.78
\notefig[0.3]{figures/p240_h.png}

\[
F=mg\sin\varphi=\frac{\mu mg}{\sqrt{1+\mu^2}}
\]
%%%END
%%%BLOCK id=240-08 ch=04 sec=圆周运动·圆锥摆 kind=knowledge alt=none cont=none y=0.80-0.95
\textbf{圆锥摆}

%FIG p240_i 0.15 0.822 0.275 0.91
\notefig[0.2]{figures/p240_i.png}

多体摆　顶点不重合，全在同一高度

\[
\left\{\begin{array}{l}
T\cos\theta=m\omega^2R\\
T\sin\theta=mg
\end{array}\right.
\qquad
\omega=\sqrt{\frac{g}{h}}
\]
%%%END
%%%BLOCK id=240-09 ch=04 sec=圆周运动·临界问题 kind=method alt=none cont=none y=0.78-0.95
常规方法

两绳上都有拉力。

%FIG p240_j 0.685 0.843 0.828 0.945
\notefig[0.25]{figures/p240_j.png}

\[
\omega_{\max}=\sqrt{\dfrac{g}{\dfrac{\sqrt3}{3}\,r}}
\qquad
\omega_{\min}=\sqrt{\dfrac{g}{\sqrt3\,r}}
\]
%%%END
%%%BLOCK id=240-10 ch=04 sec=圆周运动·临界问题 kind=note alt=02 cont=none y=0.93-0.96
结合$f$的两向性考查
%%%END
%%%PAGEEND 240
